\chapter{Zero-Day Options, Weeklies and Retail Options Flow}\label{ch:m1:zero-day-options-and-retail-flow}

At two in the afternoon the most actively traded option on the most
important equity index in the world has two hours to live. It was listed a
few weeks ago, traded little until this morning, and by the close it will be
worth either nothing or the distance between the index and its strike. In
2025 options expiring the same day averaged 2.3 million contracts a day in
the S\&P 500 index options, 59\% of all their volume. An option that close to
expiry is a different instrument from the three-month contract of
\cref{ch:m1:volatility-first-contact}: almost insensitive to \omterm{def:m1:volatility-first-contact:iv}{implied
volatility}, extremely sensitive to the path of the underlying in the next
minutes, and hedged, by whoever is on the other side, with a quantity of
shares that can swing from nothing to everything on a move of half a
percent. This chapter measures that swing.

\section{From monthlies to dailies}

\begin{definition}[Weekly and zero-day options]\label{def:m1:zero-day-options-and-retail-flow:zero}
A \emph{weekly option}\index{weekly option} is a listed \omterm{def:m1:option-contracts-and-exchanges:series}{option series} with an
expiry other than the standard monthly one, listed a few weeks before it
expires. A \emph{zero-day option}\index{zero-day option} (0DTE, zero days to
expiry) is any option on the day it expires: not a separate product, but the
last day in the life of a weekly, a monthly or a quarterly series.
\end{definition}

\begin{dated}{2026-09}{How an expiry every day came about}\label{dat:m1:zero-day-options-and-retail-flow:dailies}
Weekly S\&P 500 index options at first expired on Fridays only; Monday and
Wednesday expiries followed. In 2022 the exchange added Tuesday expiries, from 18 April, and
Thursday expiries, from 11 May: since then one series has expired every
trading day. In 2025 same-day options averaged 2.3 million contracts a day,
59\% of the index options' volume, in a US options market that traded 15.2
billion contracts in the year, 26\% more than in 2024.
\end{dated}

The appeal is arithmetic. An \omterm{def:m1:volatility-first-contact:atm}{at-the-money} option with a day to run costs
about $0.8\,\sigma\sqrt{T}$ of the underlying, some 0.8\% for an index at 16\%
volatility: a small premium for a position that doubles or vanishes by the
close. Buyers use it for an event (a data release at 08:30, a central-bank
decision at 14:00) or as a cheap intraday stop; sellers collect a premium that
decays to zero in hours, systematically, in size.
\Cref{fig:m1:zero-day-options-and-retail-flow:straddle} shows the buyer's side.

\begin{omfigure}
\begin{tikzpicture}
\begin{axis}[width=10cm,height=4.6cm,
  xlabel={index move by the close (\%)},ylabel={P\&L (\% of index)},
  tick label style={font=\small},label style={font=\small},
  xmin=-2.5,xmax=2.5,ymin=-1.4,ymax=1.9,grid=major,grid style={gray!25},table/col sep=comma]
\addplot[omDef,very thick] table[x=move_pct,y=pnl_pct]{figdata/markets-1/26-zero-day-options-and-retail-flow/straddle.csv};
\draw[thin] (axis cs:-2.5,0) -- (axis cs:2.5,0);
\node[font=\footnotesize,anchor=north] at (axis cs:0,-0.88) {premium 0.80\%};
\end{axis}
\end{tikzpicture}

\omcaption{Buying the \omterm{def:m1:volatility-first-contact:atm}{at-the-money} call and put that expire today, at 16\%
\omterm{def:m1:volatility-first-contact:iv}{implied volatility}, held to the close. The position pays if the index ends
more than 0.80\% from where it started, which at that volatility happens on
about two days in five. Data: the tutorial.}\label{fig:m1:zero-day-options-and-retail-flow:straddle}
\end{omfigure}

\section{Gamma near expiry}

\begin{definition}[Delta hedging and gamma]\label{def:m1:zero-day-options-and-retail-flow:gamma}
The \emph{delta} of an option is the change of its value per unit change of
the underlying; \emph{delta hedging}\index{delta hedging} is holding the
opposite quantity of the underlying so that the combined position is
insensitive to small moves. \emph{Gamma}\index{gamma} is the change of delta
per unit change of the underlying: the rate at which the hedge must be
adjusted. A position that is long options is long gamma: its hedger sells as
the price rises and buys as it falls. A short-option position's hedger does
the opposite.
\end{definition}

\begin{proposition}[At-the-money gamma grows like $1/\sqrt{T}$]\label{prop:m1:zero-day-options-and-retail-flow:gamma}
In the Black--Scholes model with zero rates the \omterm{def:m1:zero-day-options-and-retail-flow:gamma}{gamma} of an option with
strike $K$ is $\Gamma = \varphi(d_1)/(S\sigma\sqrt{T})$, with $d_1 =
\ln(S/K)/(\sigma\sqrt T) + \sigma\sqrt T/2$ and $\varphi$ the normal density. At
the money, $\Gamma \approx 0.4/(S\sigma\sqrt{T})$: halving the time to expiry
multiplies \omterm{def:m1:zero-day-options-and-retail-flow:gamma}{gamma} by $\sqrt2$, and \omterm{def:m1:zero-day-options-and-retail-flow:gamma}{gamma} is concentrated within about one
$\sigma\sqrt T$ of the strike.
\end{proposition}

\begin{omfigure}
\begin{tikzpicture}
\begin{axis}[width=10.5cm,height=4.6cm,xmode=log,x dir=reverse,
  xlabel={trading minutes to expiry (log scale)},ylabel={delta points per 1\%},
  tick label style={font=\small},label style={font=\small},
  xmin=4,xmax=10000,ymin=0,ymax=380,grid=major,grid style={gray!25},
  xtick={8190,1950,390,60,15,5},xticklabels={21 days,5 days,1 day,60,15,5},
  table/col sep=comma]
\addplot[omDef,very thick,mark=*,mark size=1.4pt] table[x=minutes,y=delta_change_pct]{figdata/markets-1/26-zero-day-options-and-retail-flow/atm_gamma.csv};
\draw[thin,dashed] (axis cs:4,100) -- (axis cs:10000,100);
\node[font=\footnotesize,anchor=south west] at (axis cs:9000,102) {100: the whole position};
\end{axis}
\end{tikzpicture}

\omcaption{\omterm{def:m1:volatility-first-contact:atm}{At-the-money} \omterm{def:m1:zero-day-options-and-retail-flow:gamma}{gamma} at 16\% volatility, as the change of delta, in
delta points, that a linear extrapolation gives for a 1\% move. It passes 100,
the whole position, with about an hour to go: from then on a 1\% move is more
than the option's entire remaining uncertainty. Data: the chapter's
build.}\label{fig:m1:zero-day-options-and-retail-flow:atm}
\end{omfigure}

\begin{omfigure}
\begin{tikzpicture}
\begin{axis}[width=10.5cm,height=4.8cm,
  xlabel={underlying (strike 100)},ylabel={delta change per 1\% move},
  tick label style={font=\small},label style={font=\small},
  xmin=94,xmax=106,ymin=0,ymax=1.1,grid=major,grid style={gray!25},
  legend columns=3,legend style={font=\footnotesize,at={(0.5,-0.32)},anchor=north,draw=none,fill=none,/tikz/every even column/.append style={column sep=8pt}},
  table/col sep=comma]
\addplot[omExo,very thick,densely dotted] table[x=spot,y expr=\thisrow{month}/100]{figdata/markets-1/26-zero-day-options-and-retail-flow/gamma_profile.csv};
\addplot[omDef,very thick,dashed] table[x=spot,y expr=\thisrow{day}/100]{figdata/markets-1/26-zero-day-options-and-retail-flow/gamma_profile.csv};
\addplot[omThm,very thick] table[x=spot,y expr=\thisrow{hour}/100]{figdata/markets-1/26-zero-day-options-and-retail-flow/gamma_profile.csv};
\legend{one month to expiry,one day,one hour}
\end{axis}
\end{tikzpicture}

\omcaption{\omterm{def:m1:zero-day-options-and-retail-flow:gamma}{Gamma} of a 100-strike option at 16\% volatility, expressed as the
change of delta for a 1\% move of the underlying (linear extrapolation: a
delta cannot change by more than one). With a month to go the hedge changes
by 0.09 per 1\%, whatever the price; with an hour to go, by the whole position
if the price is at the strike and by nothing if it is 1\% away. Data: the
chapter's build.}\label{fig:m1:zero-day-options-and-retail-flow:profile}
\end{omfigure}

Near expiry \omterm{def:m1:zero-day-options-and-retail-flow:gamma}{gamma} is all or nothing. The hedger of an option struck at the
current price must trade its entire notional back and forth as the
underlying crosses the strike; the hedger of an option struck 1\% away holds a
fixed position and waits. The \omterm{def:m1:option-contracts-and-exchanges:pin}{pin risk} of
\cref{def:m1:option-contracts-and-exchanges:pin} is the limit of this picture
at the closing bell.

\section{Dealer hedging and the underlying}

\begin{definition}[Dealer gamma]\label{def:m1:zero-day-options-and-retail-flow:dealer}
\emph{Dealer gamma}\index{dealer gamma} is the aggregate \omterm{def:m1:zero-day-options-and-retail-flow:gamma}{gamma} of the
positions held by \omterm{def:m1:options-market-structure:mm}{options market makers}, who hedge, as opposed to their
customers, most of whom do not. When dealers are long \omterm{def:m1:zero-day-options-and-retail-flow:gamma}{gamma} their hedging
sells rises and buys falls and tends to damp moves in the underlying; when
they are short \omterm{def:m1:zero-day-options-and-retail-flow:gamma}{gamma} it buys rises and sells falls and tends to amplify
them.
\end{definition}

\begin{omfigure}
\begin{tikzpicture}
\begin{axis}[width=10.5cm,height=4.8cm,
  xlabel={hedgers' trade as a fraction of the previous move},ylabel={volatility ratio},
  tick label style={font=\small},label style={font=\small},
  xmin=-0.65,xmax=0.65,ymin=0.4,ymax=2.6,grid=major,grid style={gray!25},table/col sep=comma]
\addplot[omDef,very thick,mark=*,mark size=1.4pt] table[x=feedback,y=vol_ratio]{figdata/markets-1/26-zero-day-options-and-retail-flow/feedback.csv};
\draw[thin] (axis cs:0,0.4) -- (axis cs:0,2.6);
\node[font=\footnotesize,anchor=east] at (axis cs:-0.08,2.2) {hedgers long gamma};
\node[font=\footnotesize,anchor=west] at (axis cs:0.08,0.7) {hedgers short gamma};
\end{axis}
\end{tikzpicture}

\omcaption{A toy market in which each five-minute return is news plus a
fraction of the previous return contributed by hedgers' trades; the vertical
axis is \omterm{def:m1:volatility-first-contact:realised}{realised volatility} over the volatility of the news alone. Feedback of
$-0.3$ removes a quarter of the day's volatility; $+0.3$ adds 40\%. The effect
is asymmetric: amplification compounds. Data: the tutorial's
simulation.}\label{fig:m1:zero-day-options-and-retail-flow:feedback}
\end{omfigure}

The mechanism is not in dispute; its size and sign on a given day are. Three
cautions before any ``\omterm{def:m1:zero-day-options-and-retail-flow:gamma}{gamma} exposure'' number is believed. First, \omterm{def:m1:futures-contracts-and-exchanges:oi}{open
interest} does not say who is long: the usual estimate \emph{assumes} that
customers buy puts and sell calls, so that dealers are short the puts and
long the calls, and its sign is that assumption. Second, same-day options
barely appear in \omterm{def:m1:futures-contracts-and-exchanges:oi}{open interest} at all, which is published once a day from the
previous night's positions: volume that opens and closes within the session
leaves no trace. Third, customer flow in same-day options is often two-sided,
buyers and sellers of the same strikes partly offsetting each other, so that the
dealers' \emph{net} position can be a small fraction of the volume. A desk that
wants this number must estimate it from signed trades, not from \omterm{def:m1:futures-contracts-and-exchanges:oi}{open
interest}.

\section{Retail options flow}

Individual investors are a large share of options volume, above all in
short-dated options on a handful of indices, funds and large shares. Their
orders reach the market through the \omterm{def:m1:retail-flow-and-wholesaling:wholesaler}{wholesalers} and auctions of
\cref{ch:m1:options-market-structure}, and their characteristics are those of
retail equity flow (\cref{ch:m1:retail-flow-and-wholesaling}) with the
leverage turned up: small size, mostly opening purchases of out-of-the-money
calls and puts held for hours or days, concentration in the nearest expiry,
and little information about the next few minutes. For a \omterm{def:m1:what-a-trading-firm-does:market-maker}{market maker} this
flow is valuable for the same reason and risky for a new one: it is
correlated. When a crowd buys the same calls on the same share, the makers
who sold them are short \omterm{def:m1:zero-day-options-and-retail-flow:gamma}{gamma} together and buy the share together as it
rises, the dynamic of
\cref{prop:m1:stock-loan-and-short-selling:cascade} with options in place of
short covering.

\section{Tutorial: gamma into the close}\label{tut:m1:zero-day-options-and-retail-flow:gamma}

\begin{tutorial}
\textbf{Goal.} Compute delta and \omterm{def:m1:zero-day-options-and-retail-flow:gamma}{gamma} in trading time, watch \omterm{def:m1:volatility-first-contact:atm}{at-the-money}
\omterm{def:m1:zero-day-options-and-retail-flow:gamma}{gamma} grow as expiry approaches, translate a position into shares to trade
per 1\% move, and simulate the feedback of hedging on the underlying.
\textbf{End state:} the four data figures of this chapter.
\begin{tutsteps}
\item \textbf{\omterm{def:m1:zero-day-options-and-retail-flow:gamma}{Gamma}.} One formula, the same for calls and puts.
\omcode{code/firm/gex/firm_gex.py}{30}{36}{Black--Scholes gamma with zero rates.}\label{lst:m1:zero-day-options-and-retail-flow:gamma}
\item \textbf{From contracts to shares.} Linear, for small moves; and exact,
by revaluing the delta.
\omcode{code/firm/gex/firm_gex.py}{53}{61}{Shares per 1\% move, and the exact hedge trade for a given move.}\label{lst:m1:zero-day-options-and-retail-flow:shares}
\item \textbf{The sign assumption}, made explicit so that it cannot be
forgotten.
\omcode{code/firm/gex/firm_gex.py}{64}{69}{Open interest seen from the dealers' side, under a stated assumption.}\label{lst:m1:zero-day-options-and-retail-flow:dealer}
\item \textbf{Feedback.} An autoregressive return whose coefficient is the
hedgers' trade per unit of move.
\end{tutsteps}
\textbf{What to change next.} Make the feedback coefficient depend on the
distance between the price and a strike with large \omterm{def:m1:futures-contracts-and-exchanges:oi}{open interest}, positive
on one side of a threshold and negative on the other, and look for pinning.
\end{tutorial}

\section{Build: the gamma-exposure estimator}\label{bld:m1:zero-day-options-and-retail-flow:gex}

\begin{build}
\bfield{Purpose} The miniature firm's \omterm{def:m1:options-market-structure:mm}{options market maker} must know, at any
moment, how many shares its own book will force it to trade for the next 1\%
move; its research side wants the same estimate for the market, with the
assumptions labelled.
\bfield{Interface} \texttt{delta(spot, strike, years, vol, right)};
\texttt{gamma(spot, strike, years, vol)}; \texttt{Holding(strike, right,
years, vol, contracts, multiplier)}; \texttt{net\_delta\_shares(book, spot)};
\texttt{gamma\_shares\_per\_pct(book, spot)}; \texttt{hedge\_trade(book,
spot\_before, spot\_after)};
\texttt{dealer\_book\_from\_open\_interest(open\_interest, years, vol,
dealer\_side)}.
\bfield{Rules} Time in trading minutes (252 days of 390 minutes), since
calendar time misprices an option with two hours to live. Expired options have
a delta of one or zero and no \omterm{def:m1:zero-day-options-and-retail-flow:gamma}{gamma}. The linear number is reported with the
exact hedge for a stated move beside it, because near expiry the two differ
by a factor of two or more. The market-wide estimate takes its sign
convention as an explicit argument and has no default.
\bfield{Acceptance tests} \texttt{code/firm/gex/tests/}: \omterm{def:m1:zero-day-options-and-retail-flow:gamma}{gamma} as the
numerical derivative of delta, equal for calls and puts; the $1/\sqrt T$ law;
the direction of hedging for long and short \omterm{def:m1:zero-day-options-and-retail-flow:gamma}{gamma}; expired options; two sign
conventions giving opposite answers on the same \omterm{def:m1:futures-contracts-and-exchanges:oi}{open interest}.
\bfield{Stretch} Signed-trade classification from the options tape to
replace the assumption, by customer and market-maker capacity codes where
the feed provides them.
\end{build}

\begin{omsources}
\item Cboe, \emph{The State of the Options Industry: 2025}; Cboe press
release, ``Cboe to Add Tuesday and Thursday Expirations for SPX Weeklys
Options'', 13 April 2022.
\end{omsources}

\section{Exercises}

\begin{exercise}[$\star$]\label{exo:m1:zero-day-options-and-retail-flow:1}
Using $0.8\,\sigma\sqrt{T}$, price the \omterm{def:m1:volatility-first-contact:atm}{at-the-money} straddle expiring today on
an index at 16\% volatility, as a percentage of the index. Same at 32\%.
\end{exercise}

\begin{exercise}[$\star$]\label{exo:m1:zero-day-options-and-retail-flow:2}
A desk is long 500 \omterm{def:m1:volatility-first-contact:atm}{at-the-money} calls (multiplier 100) on a share at 100 with
one day to expiry and 16\% volatility. The call's delta is 0.50 and its \omterm{def:m1:zero-day-options-and-retail-flow:gamma}{gamma}
0.396 per dollar. Give the hedge, and the linear estimate of the shares to
trade if the share rises 1\%. In which direction?
\end{exercise}

\begin{exercise}[$\star$]\label{exo:m1:zero-day-options-and-retail-flow:3}
From \cref{prop:m1:zero-day-options-and-retail-flow:gamma}, by what factor
does \omterm{def:m1:volatility-first-contact:atm}{at-the-money} \omterm{def:m1:zero-day-options-and-retail-flow:gamma}{gamma} grow between one month (21 days) and one day to
expiry? Between one day (390 minutes) and 15 minutes?
\end{exercise}

\begin{exercise}[$\star\star$]\label{exo:m1:zero-day-options-and-retail-flow:4}
For the position of exercise 2 the exact delta of the call at 101 is 0.84.
Give the exact hedge trade for the 1\% rise and compare it with the linear
estimate. Why does the linear estimate overshoot?
\end{exercise}

\begin{exercise}[$\star\star$]\label{exo:m1:zero-day-options-and-retail-flow:5}
Same-day index options trade 2.3 million contracts a day with a multiplier of
\$100 and the index near 6\,000. Give the notional traded. A commentator
concludes that ``1.4 trillion dollars of hedging hits the market every day''.
Give two reasons why that is wrong by orders of magnitude.
\end{exercise}

\begin{exercise}[$\star\star$]\label{exo:m1:zero-day-options-and-retail-flow:6}
An \omterm{def:m1:options-market-structure:mm}{options market maker}'s book in one share is short \omterm{def:m1:zero-day-options-and-retail-flow:gamma}{gamma} after a day of
retail call buying. Describe its hedging trades if the share rises 3\% in the
last hour, then falls back 3\% the next morning. What has it earned or lost
from hedging alone, qualitatively, and what was it paid for that?
\end{exercise}

\begin{exercise}[$\star\star\star$]\label{exo:m1:zero-day-options-and-retail-flow:7}
\emph{Coding.} With \texttt{simulate\_with\_hedgers(4000, 78, 0.001, f, 26)}
report the volatility ratio for $f = -0.3$, 0 and $+0.3$. Derive the ratio for
a long horizon analytically from the autoregression and compare.
\end{exercise}

\begin{exercise}[$\star\star\star$]\label{exo:m1:zero-day-options-and-retail-flow:8}
\emph{Find the flaw.} A service publishes each morning: ``\omterm{def:m1:zero-day-options-and-retail-flow:dealer}{Dealer gamma} is
$+\$4$ billion per 1\%: volatility will be suppressed today.'' It computes the
number from last night's \omterm{def:m1:futures-contracts-and-exchanges:oi}{open interest} assuming dealers are long every call
and short every put. Name three reasons why the number may have the wrong
size or sign for same-day options.
\end{exercise}

\section{Problem: The Last Hour}

\begin{problem}[{Weekend problem --- two million shares per percent}]\label{pb:m1:zero-day-options-and-retail-flow:1}
A dealer is short 20\,000 calls (multiplier 100) on an index fund trading at
600, struck at 600 and expiring at today's close. Volatility is 16\%; time is
measured in trading minutes, 390 to a day and 252 days to a year. The dealer
is delta hedged.

\medskip\noindent\textbf{Part I --- The hedge now.}
\begin{enumerate}
\item Give the notional of the position in shares and in dollars.
\item With 60 minutes left the call's delta is about 0.50. Give the hedge.
\item Give the one-standard-deviation move of the fund over the remaining
hour, in percent.
\item The \omterm{def:m1:zero-day-options-and-retail-flow:gamma}{gamma} is 0.168 per dollar. Give the linear number of shares to
trade per 1\% move, with its direction.
\item Compare with the notional. What does the comparison say about the
linear number?
\end{enumerate}

\medskip\noindent\textbf{Part II --- Exact trades.}
The build gives the exact hedge trades: for $+0.25\%$, buy 472\,000 shares; for
$+1\%$, buy 987\,000.
\begin{enumerate}[resume]
\item Give the delta of the call after each move.
\item Why is the second trade only about twice the first, for a move four
times larger?
\item At the open (390 minutes left) the linear number was 792\,000 shares
per 1\%. With 15 minutes left it is 4.04 million. Verify the ratio to the
60-minute figure with the $1/\sqrt T$ law.
\item The fund trades 60 million shares a day. Express the 60-minute linear
number as a fraction of an hour's average volume.
\end{enumerate}

\medskip\noindent\textbf{Part III --- The close.}
The fund is at 600.30 with five minutes left.
\begin{enumerate}[resume]
\item Describe the dealer's hedging over the next five minutes if the price
oscillates between 599.70 and 600.30.
\item What does each oscillation cost the dealer?
\item What was the dealer paid for bearing this, and by whom?
\item The options are physically settled. What else must the dealer think
about at 16:00?
\item Had the dealer been \emph{long} the 20\,000 calls, describe its trades
and their effect on the fund's price near 600.
\end{enumerate}

\medskip\noindent\textbf{Part IV --- Judgement.}
\begin{enumerate}[resume]
\item Are same-day options a threat to market stability? Give the mechanism
by which they could be, and the fact about customer flow that limits it.
\item Why is \omterm{def:m1:futures-contracts-and-exchanges:oi}{open interest} a poor guide to same-day positioning?
\item A seller of same-day straddles earns the premium on three days in
five. What does the P\&L distribution look like, and what sizes the
position?
\item Why do exchanges like daily expiries?
\item State the \emph{named result}: the shares to trade per 1\% move with an
hour left, by the linear estimate.
\item In one sentence: what is different about an option on its last day?
\end{enumerate}
\end{problem}

\section{Interview questions}

\begin{interviewq}[$\star$ \iqroles{trader, researcher}]\label{iq:m1:zero-day-options-and-retail-flow:1}
What is \omterm{def:m1:zero-day-options-and-retail-flow:gamma}{gamma}, and what does being long or short \omterm{def:m1:zero-day-options-and-retail-flow:gamma}{gamma} mean for how you
hedge?
\end{interviewq}

\begin{interviewq}[$\star$ \iqroles{trader, researcher}]\label{iq:m1:zero-day-options-and-retail-flow:2}
How does the \omterm{def:m1:zero-day-options-and-retail-flow:gamma}{gamma} of an \omterm{def:m1:volatility-first-contact:atm}{at-the-money} option behave as expiry approaches?
And of an out-of-the-money one?
\end{interviewq}

\begin{interviewq}[$\star\star$ \iqroles{trader, researcher}]\label{iq:m1:zero-day-options-and-retail-flow:3}
Price a same-day \omterm{def:m1:volatility-first-contact:atm}{at-the-money} straddle in your head. What is its break-even?
\end{interviewq}

\begin{interviewq}[$\star\star$ \iqroles{researcher, mle}]\label{iq:m1:zero-day-options-and-retail-flow:4}
How would you estimate dealers' \omterm{def:m1:zero-day-options-and-retail-flow:gamma}{gamma} positioning, and how would you test
whether your estimate predicts anything?
\end{interviewq}

\begin{interviewq}[$\star\star$ \iqroles{trader}]\label{iq:m1:zero-day-options-and-retail-flow:5}
You are short \omterm{def:m1:zero-day-options-and-retail-flow:gamma}{gamma} into the close with the underlying at your strike. What
are your choices?
\end{interviewq}

\begin{interviewq}[$\star\star\star$ \iqroles{researcher, trader}]\label{iq:m1:zero-day-options-and-retail-flow:6}
Do same-day options increase or decrease the volatility of the index? Build
the argument, then say what data would settle it.
\end{interviewq}
